\documentclass{article}
\usepackage[a4paper, margin=1in]{geometry} % Adjust margins
\usepackage{amsmath, amssymb, amsthm}
\usepackage{tikz}
\usepackage{xcolor}
\usepackage{graphicx}
\usepackage{listings}
\usepackage{fancyhdr}
\usetikzlibrary{automata, positioning}
\usepackage{helvet} % Match font to reference
\pagestyle{fancy}
\usepackage{enumitem,xcolor}


% Ensure proper header spacing
\setlength{\headheight}{15pt} % Adjust height of the header
\setlength{\headsep}{10pt} 

% Define Header/Footer
\fancyhf{}
\fancyhead[L]{\textbf{Advanced Topics in Computer Science}} % Left header
\fancyhead[R]{Chomsky Normal Form Converter} % Right header
\cfoot{\thepage} % Page number at center bottom

% Define colors for section and subsection titles
\usepackage{titlesec}
\definecolor{sectioncolor}{rgb}{0.1, 0.3, 0.7} % Dark blue
\titleformat{\section}{\large\bfseries\color{sectioncolor}}{\thesection}{1em}{}
\titleformat{\subsection}{\bfseries\color{sectioncolor}}{\thesubsection}{1em}{}

% Define custom colors for different text types
\usepackage{tcolorbox}
\tcbuselibrary{theorems}
\newcommand{\important}[1]{\textcolor{red}{#1}}
\newtcbtheorem
  []% init options
  {definition}% name
  {Definition}% title
  {%
    colback=blue!5,
    colframe=blue!35!black,
    fonttitle=\bfseries,
  }% options
  {def}% prefix
\newcommand{\note}[1]{\textcolor{green!50!black}{#1}}

% Theorem environment setup
\newtheorem{theorem}{Theorem}
\newtheorem{lemma}{Lemma}
\newtheorem{corollary}{Corollary}

\definecolor{codegray}{gray}{0.9}
\lstdefinestyle{customc}{
  backgroundcolor=\color{codegray},
  language=C,
  basicstyle=\ttfamily\footnotesize,
  keywordstyle=\color{blue},
  commentstyle=\color{green!50!black},
  stringstyle=\color{red}
}
\begin{document}

\title{\textbf{Chomsky Normal Form Converter}}
\author{Scooter Computer}
\date{}
\maketitle
\thispagestyle{fancy}

The following is a conversion of a context-free grammar (CFG) into Chomsky Normal Form (CNF). \\

This conversion has five steps:
\begin{enumerate}
\item Ensure the start variable is not on the RHS
\item Eliminate $\epsilon$-rules by first finding nullable variables
\item Eliminate unit rules
\item RHS is either one terminal or just variables
\item Ensure RHS has at most two variables
\end{enumerate}

Right-hand side will be abbreviated to RHS and left-hand side will be abbreviated to LHS.

Consider the following grammar:
\begin{align*}
    S &\to AaB  | b  | S \\
    A &\to S  | \varepsilon  | AB \\
    B &\to bbb  | ASA \\
    \end{align*}

\section{Ensure the start variable is not on the RHS}




  
    
    
  
    
    
  
    
      
    
    
  

  
    
      
    
    
  
    
    
  
    
    
  

  
    
    
  
    
      
    
    
  




  The grammar contains productions with S on the right-hand side. To eliminate these productions, simply add a new state as follows:
  \begin{align*}
    S_0 \to S
  \end{align*}


\section{Eliminate $\epsilon$-rules}
First, we need to find nullable variables in the grammar, which are any variables that can produce $\epsilon$. In this grammar,
the following variables are nullable:
\begin{align*}
    V_{\epsilon}=\{ A\}
\end{align*}

For each appearance of the nullable variable in the grammar, add a rule where that variable is replaced by $\epsilon$. This 
results in the following grammar:
\begin{align*}
    S &\to b  | S  | aB  | AaB \\
    A &\to AB  | \varepsilon  | S  | B \\
    B &\to SA  | S  | bbb  | AS  | ASA \\
    S0 &\to S \\
    \end{align*}

Finally, remove all $\epsilon$-rules to yield
\begin{align*}
    S &\to b  | S  | aB  | AaB \\
    A &\to AB  | S  | B \\
    B &\to SA  | S  | bbb  | AS  | ASA \\
    S0 &\to S \\
    \end{align*}

\section{Eliminate unit rules}


The grammar contains unit rules. Recall that unit rules are rules of the form:
\begin{align*}
    A \to B \\
    B \to C | D\\
\end{align*}
Notice that $B$ can be eliminated via $A \to C | D$. 
Essentially, we substitute all of the rules of $B$ into $A$. 
One way to visualize this is through a directed graph.

\begin{figure}[h]
    \centering
    \includegraphics[width=0.5\linewidth]{phase2/chomsky_converter/out/directed_graph.png}
    \caption{A visual representation of which variables can be reached via a series of applications of unit rules.}
    \label{fig:dir_graph}
\end{figure}

We will perform this operation for all unit rules
in the provided grammar, which yields:
\begin{align*}
    S &\to b  | aB  | AaB \\
    A &\to b  | SA  | bbb  | AS  | ASA  | aB  | AaB  | AB \\
    B &\to b  | SA  | bbb  | AS  | aB  | AaB  | ASA \\
    S0 &\to b  | aB  | AaB \\
    \end{align*}



\section{Ensure RHS is either one terminal or just variables}
If any rule contains a combination of terminals and variables, then simply replace that terminal with a variable. 
For example, the terminal $a$ would be replaced by $U_a$. This yields the following grammar:
\begin{align*}
    S &\to b  | U_aB  | AU_aB \\
    A &\to b  | SA  | AS  | ASA  | AB  | AU_aB  | U_aB  | U_bU_bU_b \\
    B &\to b  | SA  | AS  | ASA  | AU_aB  | U_aB  | U_bU_bU_b \\
    S0 &\to b  | U_aB  | AU_aB \\
    [Ua] &\to a \\
    [Ub] &\to b \\
    \end{align*}

\section{Break up long RHS rules}
Finally, we need to break up rules that are longer than two variables. The simplest approach for this is to 
go through the rules that are too long and combine variables until they are exactly two variables long. Once
complete, this yields the following grammar:
\begin{align*}
    S &\to b  | U_aB  | CB \\
    A &\to b  | SA  | AS  | AB  | U_aB  | DA  | CB  | EU_b \\
    B &\to b  | SA  | AS  | U_aB  | DA  | CB  | EU_b \\
    S0 &\to b  | U_aB  | CB \\
    [Ua] &\to a \\
    [Ub] &\to b \\
    C &\to AU_a \\
    D &\to AS \\
    E &\to U_bU_b \\
    \end{align*}

Now, notice that all rules point to either a single terminal or exactly two variables. This grammar is now in Chomsky Normal Form.
\end{document}